% B.41: independently audited by C; exact code is in deliverable/code.
\subsection{An arithmetic action term not determined by the Ross extension}

\begin{statement}{B.41}{The Ross extension does not determine arbitrary-input lifting}{Proved}
Keep the actual cubic group, lattice, and $C_0$ tangent of B.24--B.40.
Write $N_3=N/\gamma_4N$, $N_2=N/\gamma_3N$, and
$A=\gamma_3N/\gamma_4N=\operatorname{Lie}_3(V)$, without a Tate twist.
There is a continuous arithmetic action $\rho'$ on the same group
$N_3$ with the following properties.

It has the same action as the actual action $\rho$ on $N_2$ and on
$\gamma_2N/\gamma_4N$.  In particular it has the same entire
extension defining $\eta$, the same rational Ross class, and the
same integral refinement determined in B.40.  It fixes the chosen
$C_0$ inertia, preserves every peripheral conjugacy class, and has
the same compatibility with the outer geometric action of
$\mathcal A=E[3]\rtimes\mu_6$.

Nevertheless the class-two input obtained from the actual $A_0$
cusp section by the twist $\kappa_3(19)$ lifts for $\rho$ and does
not lift for $\rho'$, already locally at nineteen.  The modified
action is not asserted to be the action of an algebraic curve.
This is an insufficiency theorem for the data determined in task
B.4, not a determination of the actual image requested in B.5.

The other rational cusps' class-two origins change under this
deformation by an explicit central path cocycle.  Thus merely
requiring the existence of all peripheral packets does not remove
the deformation.  Fixing an independently known actual $A_0$
class-two origin and all its Kummer twists would remove it.
\end{statement}
\begin{proof}

\par\smallskip\noindent\textbf{An explicit integral central derivation.}\enspace
Let $W$ be the boundary lattice with basis
$(A_0,A_1,A_2,B_0,B_1,B_2,C_1,C_2)$ and with
$e_O=-\sum_{P\ne O}e_P$, where $O=C_0$.  In the following
brackets these cusp names mean their boundary vectors.  Put
\[
 b=[B_0,A_2-A_1]+[B_1,A_0-A_2]+[B_2,A_1-A_0]
       \in\Lambda^2W.                                      \tag{B.41.1}
\]
This vector is fixed by the stabilizer $\mu_6$ of $O$.
For $P\in E[3]$, let $T_P$ be the actual translation matrix on
$V$, and define a rational degree-two derivation by
\[
 D(u)=D(v)=0,\qquad D(e_P)=[e_P,T_Pb].                       \tag{B.41.2}
\]
Here $(u,v)$ is the primitive compact pair of B.24.  The nine
values in (B.41.2) sum to zero, so this is well defined on $W$.
The stabilizer condition makes it $\mathcal A$-equivariant.
These identities are exact finite identities on the nine-point
affine plane; the replayable certificate evaluates all 54
geometric matrices and all nine boundary vectors.

In the integral basis $(f_1,f_2,W)$ of B.24, write $D_W$ for
the last eight columns of $D$.  Since $(u,v)=(3I_2,U)$,
\[
 D=\bigl(-D_WU/3\mid D_W\bigr).
\]
Define
\[
 \Delta=3D+3\operatorname{ad}_b
       =\bigl(-D_WU\mid3D_W\bigr)+3\operatorname{ad}_b
       \in\operatorname{Hom}_{\mathbb Z_3}(V,A).             \tag{B.41.3}
\]
Thus $\Delta$ is integral.  It is primitive, fixes $e_O$, and
at every boundary vector satisfies
\[
 \Delta(e_P)=3[e_P,B_P],\qquad B_P=T_Pb-b.                  \tag{B.41.4}
\]
For a geometric automorphism $a\in\mathcal A$ one has
\[
 \operatorname{Lie}_3(a)\Delta-\Delta a
    =\operatorname{ad}_{3(ab-b)}\,a.                       \tag{B.41.5}
\]
Consequently the outer central derivation commutes with the
geometric action, with the displayed integral inner correction.
For elements fixing $O$, the correction is zero.

For clarity, these computations use the original integral
basis, not the rational direct sum $H\oplus W$.  The first
parameter vector in \texttt{b\_arbitrary\_action\_analysis}
has exactly the six coefficients displayed in (B.41.1).
The program \texttt{b\_arbitrary\_action\_deformation}
constructs (B.41.3) and checks (B.41.4)--(B.41.5) over
unreduced integers.  The independent certificate
\texttt{c\_B41\_exact\_audit} reconstructs the six terms
and all boundary columns without importing the first
calculation; it obtains the same integral matrices.

\par\smallskip\noindent\textbf{Its actual arithmetic coefficient module.}\enspace
Let $\epsilon=\chi_{-3}$.  The arithmetic Frobenius and
conjugation matrices from B.28 give
\[
 \operatorname{Lie}_3(T_V)\Delta=49\Delta T_V,
 \qquad
 \operatorname{Lie}_3(C_V)\Delta=-\Delta C_V.               \tag{B.41.6}
\]
Hence the line $\Lambda=\mathbb Z_3\Delta$ has the actual
coefficient action $\mathbb Z_3(2)\epsilon$.  This also
follows directly from (B.41.1): arithmetic over
$K=\mathbb Q(\zeta_3)$ scales the boundary by $\chi_3$,
complex conjugation sends $b$ to $-b$, and the rational
compact splitting is arithmetic stable.

Use the arithmetic quotient $\Gamma=P\rtimes\langle c\rangle$
of B.37, where $P$ is free pro-three on the two actual
Frobenii $\tau,\bar\tau$ above seven and $c$ interchanges
them.  A $\Lambda$-valued cocycle has a unique representative
with value zero at $c$: multiplication by two is invertible,
and $\Lambda^{c=1}=0$.  The values on the free basis are then
$a\Delta,-a\Delta$, with arbitrary $a\in\mathbb Z_3$.
No nonzero coboundary preserves the normalization at $c$.
Therefore
\[
 H^1(\Gamma,\Lambda)\simeq\mathbb Z_3.                    \tag{B.41.7}
\]
Take the cocycle $\psi(g)=b_g\Delta$ with
\[
 b_c=0,\qquad b_\tau=1,\qquad b_{\bar\tau}=-1.             \tag{B.41.8}
\]
Freeness gives a continuous cocycle: equivalently, prescribe
the two images in the pro-three semidirect product
$\Lambda\rtimes\rho_\Lambda(P)$ and apply the defining
universal property of the free pro-three group.  The
conjugation relation in (B.41.8) extends it to $\Gamma$.

Every integral map $\delta:V\to A$ defines a literal central
automorphism of $N_3$ by $n\mapsto n\,\delta(n^{\rm ab})$.
Its inverse uses $-\delta$.  It is the identity on $N_2$
and on $\gamma_2N/\gamma_4N$; this uses only that $A$ is
central and does not require a class-three exponential map.
Set
\[
 \rho'_g=\mathcal D_{\psi(g)}\circ\rho_g.                 \tag{B.41.9}
\]
The cocycle identity is exactly the action identity for
(B.41.9).  Pulling back along $G_{\mathbb Q}\to\Gamma$
gives a continuous arithmetic action unramified outside
three.  Equations (B.41.3)--(B.41.5) give all the preservation
properties stated above.  In particular the arithmetic
and geometric actions on the extension used in B.40 are
unchanged, since central automorphisms of this kind act
trivially on $\gamma_2N/\gamma_4N$.

\par\smallskip\noindent\textbf{The class-two input and its local failure.}\enspace
Choose the actual global $A_0$ cusp section with a rational
unit tangent, and twist it by $\kappa_3(19)$ in its inertia.
Its projection is a class-two input for both actions, since
their $N_2$ actions coincide.  It has a full lift for $\rho$
by its construction as an actual cuspidal section.

At the prime of $K$ above nineteen for which
$\zeta_3\equiv7$, put $\lambda=1-\zeta_3$.  The two
Kummer rows on $\tau,\bar\tau$ are $(2,0),(2,2)$, so
\[
 b\bmod3=2\kappa(\zeta_3)+2\kappa(\lambda).
\]
Here $\kappa(\zeta_3)(F_{19})=0$ and
$13^6\equiv11=7^2\pmod {19}$, giving
\[
 b_{F_{19}}\equiv1\pmod3.                                 \tag{B.41.10}
\]
Let $x$ be the lifted tame inertia generator of the actual
cuspidal section.  Its abelianization is $e_{A_0}$.  Write
\[
 \Delta(e_{A_0})=3z.
\]
The vector $z$ is integral, primitive, and lies in the pure
boundary summand $\operatorname{Lie}_3(W)$.  For example,
in the original Hall basis its coefficient at $(3,2,5)$
is one; the complete thirty-term vector is saved in the
certificate JSON.  The deformation changes the tame
Frobenius relation by $19b_{F_{19}}\Delta(e_{A_0})$.
Changing the central Frobenius lift has no effect on this
relation.  Changing the inertia lift by $a\in A$ changes
it by $(\Phi_A-19)a$.  Thus the unique rational correction
that would restore the relation is
\[
 a=-\frac{3b_{F_{19}}}{19^2-1}\,z
   =-\frac{b_{F_{19}}}{120}\,z.                            \tag{B.41.11}
\]
Indeed $\Phi_A=19^3$ on the pure boundary summand.
On all of $A_{\mathbb Q_3}$ the operator
$\Phi_A-19$ is invertible: every degree-three weight has
complex absolute value greater than nineteen.
Equation (B.41.10) and the primitive coefficient of $z$
show that (B.41.11) has valuation $-1$.  It is not integral,
so there is no local lift, and hence no global lift for
$\rho'$.  Conjugating a representative of the class-two
input cannot alter this conclusion, since any conjugating
element of $N_2$ lifts to $N_3$.

\par\smallskip\noindent\textbf{Exactly which cusp information has changed.}\enspace
If $P$ is a rational cusp, $B_P$ is an arithmetic eigenvector
of type $\mathbb Z_3(2)\epsilon$.  Consequently
\[
 \delta_P(g)=3b_gB_P\in\Lambda^2V                         \tag{B.41.12}
\]
is a genuine global one-cocycle.  After fixing the abelian origin and the chosen peripheral
generator, the class-two origin of the peripheral packet for the modified action must differ
from the actual origin by this cocycle.  To check the sign,
the added arithmetic action on $e_P$ is
$3\chi_3(g)b_g[e_P,B_P]$; conjugation by an element with
degree-two part $3b_gB_P$ adds
$3\chi_3(g)b_g[B_P,e_P]$ and cancels it.
The map $\operatorname{ad}_{e_P}:\Lambda^2V\to A$ is
injective, as is seen after extending the primitive
$e_P$ to a free basis.  Thus this degree-two change is forced.

The origin at $O$ stays fixed because $B_O=0$.  For $A_0$
the change is nonzero.  Its cocycle class is nonzero
rationally as well: after normalizing at $c$, a possible
coboundary must come from the $c$-fixed subspace, whereas
$B_{A_0}$ lies in the $c$-anti-invariant pure boundary
summand and $\Lambda^2(T_V)$ acts there by the scalar $49$.
The arithmetic and conjugation stable rational splitting
$H\oplus W$ separates this summand from the mixed and
compact summands, so a conjugator in those summands cannot
supply its missing anti-invariant component.
The modified class-two peripheral origin has a class-three
lift over $\Gamma$: its remaining central obstruction is
in $H^2(\Gamma,A)=0$, because $P$ is free and two is a unit.
Central corrections do not affect the peripheral relation.
It therefore provides a formal peripheral packet, whose
origin is (B.41.12), rather than preserving the old packet
pointwise.

If the actual $A_0$ class-two origin were independently
known and held fixed, all its Kummer twists would detect
any nonzero scalar multiple $a\psi$.  In fact
$b_{\tau\bar\tau^{-1}}=2$ and this element has cyclotomic
character one.  Chebotarev supplies primes whose Frobenius
is arbitrarily close to it in any chosen finite quotient.
For any nonzero $a$, choose such a prime $p$ with
$v_3(p-1)>1+v_3(a)$ while
$v_3(ab_{F_p})=v_3(a)$.  The same local calculation then
requires
\[
 -\frac{3ab_{F_p}}{p^2-1}z\in A,
\]
which is impossible.  This explains both the strength and
the limitation of imposing actual cusp inputs: it needs
the actual central path origin (B.41.12), not only the
existence of a packet or its abelianized centre.

\par\smallskip\noindent\textbf{The precise remaining datum.}\enspace
This example exhibits one missing cubic arithmetic-action
component, of type $H^1(\Gamma,\mathbb Z_3(2)\epsilon)$,
equivalently a component of a class-two cusp-path origin.
It does not claim that this is the only datum required
for arbitrary inputs.  The independent rational module
calculation gives two pure boundary tangential directions
of this type: the stabilizer-invariant conjugating
parameters have rank seven; the nine-cusp relation leaves
rank three, consisting of one inner $\Lambda^2H$ direction
and two $\Lambda^2W$ directions, with no $H\otimes W$
direction.  The full actual values of these directions
have not been computed here.  B.28 still evaluates every
nonempty fibre, and B.31 and B.36 still determine the
stated zero-abelian slice.  Neither is an evaluated image
law for every input in task B.5.

\end{proof}
